\section{Everything has a type}\label{sec:01-basics:01-everything-has-a-type:1}

In Lean 4, every expression (a \textbf{term}) has a \textbf{type}. A type answers the question ``what \emph{kind} of thing is this?'' --- a number, a boolean, a proof, a function from numbers to numbers. Lean checks this question \emph{before} running anything, and it never lets a term of the wrong type slip through.

\begin{lstlisting}
#check 3          -- 3 : Nat
#check -3         -- -3 : Int
#check Nat        -- Nat : Type
#eval 2 ^ 10        -- 1024 (evaluates; #check only elaborates)
\end{lstlisting}

\begin{mathreading}

\texttt{\#check e} is the judgment \(e : \tau\) from type theory. It is literally the same typing judgment used in a lecture on the simply-typed (or dependently-typed) \(\lambda\)-calculus: \(3 : \mathtt{Nat}\), \({-3} : \mathtt{Int}\), \(\mathtt{Nat} : \mathtt{Type}\). \texttt{\#eval e} instead asks Lean to reduce \(e\) to a normal form. This is the computational content of \(\beta\)-reduction: Lean follows the definitional equalities until nothing more can fire, exactly as one would normalize a term in \(\lambda\)-calculus by hand. (Both ``normal form'' and ``\(\beta\)-reduction'' receive a full treatment in \hyperref[sec:01-basics:04-terminology:1]{Chapter 1 §4} for readers meeting them for the first time; briefly, each denotes the value obtained after fully simplifying/running the expression, nothing more exotic.)

\end{mathreading}

Read \texttt{\#check 3} as ``Lean answers the question \texttt{3 : ?} with \texttt{Nat}.'' Read \texttt{\#check Nat} as ``even \texttt{Nat} itself --- the type of natural numbers --- is a term, and \emph{its} type is \texttt{Type}.'' This is the first surprising fact worth sitting with: types are not a separate kind of thing bolted onto a programming language. In Lean, a type is itself a term, and it has a type of its own (\texttt{Type}), the same way \texttt{3} has a type (\texttt{Nat}). This is what ``everything has a type'' means literally, not just for ordinary values.

\texttt{\#check e} answers the question ``what type does \texttt{e} have?'' without running \texttt{e}. \texttt{\#eval e} actually runs \texttt{e} and prints the result. These are deliberately two different commands, since they answer two different questions:

\begin{itemize}
\tightlist
\item
  \texttt{\#check} is a \textbf{static} guarantee --- it holds before any particular input is supplied, for every possible run.
\item
  \texttt{\#eval} is a \textbf{one-off fact} --- the result of running this \emph{particular} expression, right now.
\end{itemize}

\begin{progcorner}

For readers who have written Python but not Lean: \texttt{\#check e} is \emph{not} \texttt{type(e)}. Python's \texttt{type()} asks a running value what class it happens to belong to, after the fact. Lean's \texttt{\#check} is closer to what a static type checker like \texttt{mypy} does with an annotation such as \texttt{x: int = 3}: it verifies, \emph{before} anything runs, that the expression could only ever produce a value of the stated type. \texttt{\#eval e}, on the other hand, \emph{is} just \texttt{print(e)}: run it, show the result. Lean separates these two commands for the same reason \texttt{mypy} exists at all: type-checking is a static guarantee that holds for every possible input, while evaluating is a one-off fact about this particular expression. Python's \texttt{int} has no genuine analogue of \texttt{Nat}: it is signed and never checked against a ``must be non-negative'' rule, except by an explicit runtime \texttt{if} statement. \texttt{Nat}, by contrast, bakes non-negativity into the type itself, checked once, statically, for every use site. This is closer to a language with a genuine \texttt{unsigned} type (C, Rust) than to Python.

\end{progcorner}

\subsection{Why this matters: types rule things out in advance}\label{sec:01-basics:01-everything-has-a-type:2}

Here is the entire point of a type system, made concrete --- starting in Python, where the failure can actually be watched happening, before seeing how Lean rules it out instead.

\begin{lstlisting}[language=Python,style=python]
# Python: this line is perfectly legal to *write*.
def add_them(a, b):
    return a + b

add_them(3, True)    # 4 (*\textemdash{}*) Python silently treats True as 1, no error at all
add_them(3, "oops")  # TypeError: unsupported operand type(s), but only once this exact line runs
\end{lstlisting}

Nothing stops Python from \emph{writing} \texttt{add\_\allowbreak{}them(3, "oops")}. The mistake is only discovered the moment that specific line executes --- if it sits on a rarely-hit branch, it can ship for months undetected. And in the \texttt{True} case, Python does not even complain: it quietly coerces the boolean to \texttt{1} and moves on, whether or not that was the intended meaning.

Now the same shape of mistake in Lean:

\begin{lstlisting}
#check 3 + true   -- error: failed to synthesize Add Nat Bool
\end{lstlisting}

Lean refuses to even \emph{elaborate} this expression. It never runs it, never silently coerces \texttt{true} to \texttt{1}, never crashes five function calls later once the bad value finally reaches code that cannot handle it --- it simply never accepts the term in the first place, because \texttt{+}'s left argument is a \texttt{Nat} and its right argument is a \texttt{Bool}, and no rule connects the two. The check happens once, by reading the expression, with no particular inputs run at all. Lean's promise is stronger than Python's: if a term type-checks, an entire class of ``it blew up unexpectedly'' bugs is already provably impossible for that term, for \emph{every} input, not merely the ones a test suite happened to cover.

\subsection{\texorpdfstring{\texttt{Nat}, concretely}{Nat, concretely}}\label{sec:01-basics:01-everything-has-a-type:3}

\texttt{Nat} is not a built-in primitive the way \texttt{int} is in most languages. It is defined, in full, as an \textbf{inductive type}:

\[
\mathtt{Nat} ::= \mathtt{zero} \mid \mathtt{succ}\,(n : \mathtt{Nat})
\]

Read this as: ``a \texttt{Nat} is built in exactly one of two ways --- it is \texttt{zero}, or it is \texttt{succ n} for some already-built \texttt{Nat} called \texttt{n}.'' This is Peano's definition, written out exactly. So \texttt{3} is not a primitive digit; it is shorthand for \texttt{succ (succ (succ zero))}. Confirm it directly:

\begin{lstlisting}
#eval Nat.succ (Nat.succ (Nat.succ Nat.zero))  -- 3
\end{lstlisting}

Lean prints numerals for readability, but underneath, every \texttt{Nat} really is built from nothing but \texttt{zero} and \texttt{succ}, the same way every natural number in a first course in logic is built from Peano's axioms. This inductive shape is exactly what licenses proof by induction later (Chapter 4): to prove something about \emph{every} \texttt{Nat}, it suffices to prove it for \texttt{zero} and show it is preserved by \texttt{succ} --- because those are, provably, the only two ways a \texttt{Nat} can ever have been built.

\begin{quote}
\textbf{Mathematical reading (optional, for readers who already know some category theory).} It is useful to regard \texttt{Type} as (the objects of) a category: a term \texttt{x : α} is an element of \texttt{α}, a function \texttt{f : α → β} is a morphism \(\alpha \to \beta\), \texttt{fun x => x} is the identity, and \texttt{∘} is genuine categorical composition, with associativity and the identity laws holding \emph{definitionally} (Lean checks them automatically, at no extra cost). In this language, \texttt{Nat} is the \textbf{initial object} in the category of ``sets equipped with a distinguished element and an endomorphism'' (\(\mathrm{zero} \in X\) and a map \(s : X \to X\)): a map out of \texttt{Nat} is uniquely determined by where it sends \texttt{zero} and how it commutes with \texttt{succ}, which is exactly the universal property that makes structural induction valid. None of this is required to use \texttt{Nat} --- it is offered only because, once \texttt{+}/\texttt{0} are \emph{defined} on \texttt{Nat} (Chapter 4), it becomes provable (not just definitional) that \((\mathbb{N}, +, 0)\) is the free commutative monoid on one generator, a genuinely different universal property from the one above, worth not conflating with it. \hyperref[sec:01-basics:04-terminology:1]{Chapter 1 §4} fixes the vocabulary (``universal property,'' ``initial object'') used here, for any reader meeting it for the first time.
\end{quote}

\subsection{References}\label{sec:01-basics:01-everything-has-a-type:4}

Full citations in the \hyperref[sec:root:bibliography:1]{Bibliography}.

\begin{itemize}
\tightlist
\item
  Lean 4 documentation, ``Basic Types,'' and \emph{Theorem Proving in Lean 4}, Ch. 2 (\cite{LeanDocs}, \cite{TPIL4}) --- the \texttt{\#check}/\texttt{\#eval} distinction and \texttt{Nat} as an inductive type, straight from the source.
\item
  Pierce (\cite{Pierce2002}), Ch. 1 --- on what a static type system buys (ruling out whole classes of runtime failure before execution), independent of any particular language.
\item
  nLab, ``initial object'' (\cite{NLabInitial}) --- the universal-property reading of \texttt{Nat} used in the optional box above.
\end{itemize}
